\chapter{Single Fields}

\section{Electrostatics}
\subsection{Results}
The electric field intensity vector $\bE$ can be derived from the
nodal degree of freedom, the electric potential $V$ by applying a
gradient operator. This element quantity is calculated by applying
the config-file-commands

\begin{verbatim}
<storeResults>
   <elemResults type="elecFieldIntensity" region="elecst2d"/>
</storeResults>
\end{verbatim}

In the command {\bf $<$region$>$}, all areas in which the electric field has to
be calculated, must be listed. If no region  is specified, the elecric
field will be calculated in all regions where the PDE is defined. To
calculate the electric energy, one has to use the command 

\begin{verbatim}      
<storeResults>
   <elemResults type="elecEnergy" region="reg2"/>
</storeResults>
\end{verbatim}
in the config-file. The result is written to the info-file.


\subsection{Example}

{\footnotesize
\begin{verbatim}
<?xml version="1.0"?>

<cfsSimulation xmlns="http://www.cfs++.org">

  <analysis type="static"/>
  <geometry type="axi"/>
  <input format="mesh"/>
  <output format="unv"/>
  <materialData file="mat.dat"/>

  <domain>
    <region name="gndEl"  material="steel"/>
    <region name="hotEl"  material="steel"/>
    <region name="dielEl" material="silicon"/>
    <region name="airEl"  material="air"/>

    <nodes name="hot"/>
    <nodes name="gnd"/>
  </domain>

  <pdeList>
    <electrostatic>
       <region name="gndEl"/>
       <region name="hotEl"/>
       <region name="dielEl"/>
       <region name="airEl"/>

      <bcsAndLoads>
        <dirichletHom name="gnd"/>
        <dirichletInhom name="hot" value="1.0"/>
      </bcsAndLoads>

      <storeResults>
        <elemResults type="elecFieldIntensity"/>
        <elemResults type="elecEnergy"/>
      </storeResults>

    </electrostatic>
  </pdeList>

</cfsSimulation>
\end{verbatim}
}


\newpage
\section{Mechanics}
\subsection{Degree of freedoms}
In the 3D case we have {\bf ux, uy, uz} and in the 2D case
(axisymmetric as well as plane) {\bf ux, uy}.

\subsection{Subtypes}
The type of 2D-element is in general set in the top-level command

\verb!   <geometry type="plane">!

\smallskip
\noindent
For the mechanic PDE, a specialization has to be done, which has to be
performed as {\bf subtype}-attribute inside the tag for the PDE:{\bf
  $<$mechanic$>$}. Plane stress models are simulated by

\verb!   <mechanic subtype="planeStress">!

\noindent
and plane strain can be defined by

\verb!   <mechanic subtype="planeStrain">!

\bigskip
\noindent
{\em Notice:} Standard 2D subtype is {\sf planeStrain} and a mixture
of element subtypes within one analysis is not allowed. 

\subsection{Nodal Results}
Possible nodal results are

\begin{verbatim}
   <nodeResults type="mechDisplacement" region="bar"/>
   <nodeResults type="mechVelocity" region="bar"/>
   <nodeResults type="mechAcceleration" region="bar"/>
\end{verbatim}

\subsection{Element Results}
The command for calculating the mechanical stresses is
\begin{verbatim}      
<storeResults>
   <elemResults type="mechStress" region="bar"/>
</storeResults>
\end{verbatim}

\subsection{Example}

{\footnotesize 
\begin{verbatim}
<?xml version="1.0"?>

<cfsSimulation xmlns="http://www.cfs++.org">

  <analysis type="static"/>
  <geometry type="plane"/>
  <input format="mesh"/>
  <output format="unv"/>
  <materialData file="mat.dat"/>

  <domain>
    <region name="cantilev"  material="steel"/>

    <nodes name="leftN"/>
    <nodes name="rightN"/>
  </domain>

  <pdeList>
    <mechanic subtype="planeStrain">
       <region name="cantilev"/>

      <bcsAndLoads>
        <dirichletHom name="leftN" dof="ux"/>
        <dirichletHom name="leftN" dof="uy"/>
        <dirichletHom name="rightN" dof="ux"/>

        <load         name="rightN" dof="uy" value="-1.0"/>
      </bcsAndLoads>

    </mechanic>
  </pdeList>

</cfsSimulation>
\end{verbatim}
}


\newpage
\section{Acoustics}

\subsection{General Setup}
The \xml{acoustic} element is appended by the following child elements
\begin{itemize}
\item \xml{region}
\item \xml{timeSteppingFormulation} to date effective stiffness formulation 
is default(optional)
\item \xml{bcsAndLoads} (optional)
\item \xml{nonLinear}\footnote{Other Nonlinearities than nonlinear acoustic
 PDE, which can be chosen by an attribute of \xml{region} element. This is
 currently not implemented.}
\item \xml{storeResults} (optional)
\item \xml{solver}(optional)
\end{itemize}
and the following attributes
\begin{itemize}
\item \xml{formulation} i.e. acoustic pressure or acoustic potential PDE 
formulation
\item \xml{subType} for flownoise
\end{itemize}

\subsection{Damping and nonlinear acoustic simulation}
Damping and nonlinear acoustic PDE formulations can be chosen for each 
region individually. Damping is a child element of \xml{region},
type is either no damping, rayleigh damping (which is not physical), 
thermoviscous damping and fractional damping according to a 
frequency power law.
\xml{nonLinear} comes as attribute to region. It can take the value 'no', 
'westervelt', or 'kuznetsov'.

\subsubsection{Nodal Results}
The config-file command for storing (only) the acoustic potential in the result
file is

\begin{verbatim}      
<storeResults>
   <nodeResults type="acouPotential" region="fluid1"/>
</storeResults>
\end{verbatim}

%\subsubsection{Element Velocity $\{v\}$}

%\begin{print}
%  The velocity can be calculated as an element result by applying an
%  {\em attgroup} command.
%\end{print}

%\begin{verbatim}
%attgroup,1,calc,calc_gvp
%\end{verbatim}


\subsection{Absorbing Boundary Conditions (ABC)} 
In the config file, the command

\begin{verbatim}    
<bcsAndLoads>  
   <absorbingBCs name="abc"/>
</bcsAndLoads>
\end{verbatim}
      
is needed, which defines the boundaries at which the ABCs are working.
It should be clear, that the dimension of the boundary elements are
one order lower compared to the dimension of the "domain" elements.
Therefore, the boundary elements have to be listed with the config-file
command

\begin{verbatim}
<domain>
   <elements name="abc"/>
</domain>
\end{verbatim}

\subsection{Formulation}
For linear acoustics, the PDE in pressure and acoustic potential has the same 
form. Allthough default is acoustic potential. If acoustics is coupled to 
other PDEs, acoustic potential is obligatory. Westervelt's nonlinear PDE is 
in pressure, Kuznetsov's nonlinear PDE is in acoustic potential.


\subsubsection{Example}

{\footnotesize 
\begin{verbatim}
<?xml version="1.0" encoding="ISO-8859-1"?>
<cfsSimulation xmlns="http://www.cfs++.org">
  <analysis type="transient"/>
  <geometry type="plane"/>
  <domain>
    <region material="softTissue" name="reg1"/>
    <region material="bone" name="reg2"/>
    <nodes name="vp-restr"/>
  </domain>
  <pdeList>
    <acoustic formulation="acouPotential">
      <region name="reg1" nonLinear="no">
        <damping type="no"/>
      </region>
      <region name="reg2" nonLinear="no"/>
      <bcsAndLoads>
        <dirichletInhom dynamics="sin1M40cos.dat" name="vp-restr" value="1.0"/>
      </bcsAndLoads>
      <storeResults>
        <nodeHistory saveNodes="historyPoints" type="acouPotential"/>
      </storeResults>
    </acoustic>
  </pdeList>
  <transient>
    <numSteps>    600     </numSteps>
    <firstDt>     2.5e-08 </firstDt>
    <stepSaveBeg>   1     </stepSaveBeg>
    <stepSaveEnd> 600     </stepSaveEnd>
    <stepSaveInc>   1     </stepSaveInc>
    <timeDataFile name="sin1M40cos.dat"/>
  </transient>
</cfsSimulation>
\end{verbatim}
}


\newpage
\section{Magnetics}
\subsubsection{Element Results}
In the case of magnetic field calculation, many element results can be
deduced by a postprocessing step:

\begin{itemize}
\item Magnetic induction (magnetic flux density) $\bB$
    \newline Derived from the magnetic vector potential $\bA$ by
    applying the curl operator.
  
\begin{verbatim}
<storeResults>
   <elemResults type="magFluxDensity" region="yoke"/>
</storeResults>
\end{verbatim}
  
\item Eddy current density $\bJ_{\rm e}$ 
    \newline Currents induced by a magnetic field in a conductive
    material. 

\begin{verbatim}
<storeResults>
   <elemResults type="magEddyCurrent" region="yoke"/>
</storeResults>
\end{verbatim}
  
  
\item Magnetic energy $W$ 
    \newline Measure for the energy stored in the magnetic assembly.
  
\begin{verbatim}
<storeResults>
   <elemResults type="magEnergy" region="yoke"/>
</storeResults>
\end{verbatim}

\item  Inductivity $L$ 

\begin{verbatim}
<saveFileL>   lFile.dat   </saveFileL>
\end{verbatim}

The inductivity $L$ of the magnetic assembly, as well as
the magnetic flux $\Phi$ are calculated and stored in a file with
the name "lFile.dat"  given in column oriented form:

{\em time inductivity magFlux}.  
  

%\item Magnetic force $F_{\rm V}$, {\em calc\_fv}, $<$emag$>$-file
%  \newline Used for calculating the magnetic force $F_{\rm V}$, due to
%  a current loaded coil within a magnetic field (Lorentz force)..

\end{itemize}

\subsection{Coil Definition}

The list of coils has to be defined inside the tags

\begin{verbatim}
<coils>

</coils>
\end{verbatim}

\subsubsection{Coil Types}

In our simulation, we distinguish different coil types:
\begin{itemize}
\item Measurement coils
  \begin{itemize}
  \item 2D: \verb! <measurementCoil2d name="coil1">!
  \item 3D: \verb! <measurementCoil3d name="coil2">!
  \end{itemize}
\item Volage driven coils 
  \begin{itemize}
  \item 2D: \verb! <voltageCoil2d name="coil3">!
  \item 3D: \verb! <voltageCoil3d name="coil4">!
  \end{itemize}
\item Current driven coils 
  \begin{itemize}
  \item 2D: \verb! <currentCoil2d name="coil5">!
  \item 3D: \verb! <currentCoil3d name="coil6">!
  \end{itemize}
\end{itemize}


There are several parameters which must be defined for a coil. Every
coil needs an attribute {\bf name="xyz"}, standing nearby the
definition of the coil:

\verb! <voltageCoil2d name="coil3">!

\noindent
Further on, the cross section of a winding has to be
specified for every coil:

\verb!<windingCrossSection> 1e-6 </windingCrossSection>!

\noindent
The number of windings $N$ results from the cross section of the wire. It
is calculated by the overall cross section of the coil $A_c$ devided
by the cross section of a winding $A_w$

\begin{equation}
  \label{eq:N}
  N=\frac{A_c}{A_w}
\end{equation}

\noindent
Depending on the coil type, several additional tags are needed. They
are described in the next sections.


\subsubsection{Measurement Coils}

\begin{itemize}
\item Coil voltage $U$ respectively coil current $I$

\begin{verbatim}
<saveFileU>  uiFile.dat  </saveFileU>
\end{verbatim}

  With this tag the voltage respectively the current of a
  special coil element can be calculated. It's stored in a file
  named {\em resultFile} in the column oriented form:

  {\em time voltage}.

\end{itemize}

\subsubsection{Current- and Voltage Driven Coils}
In this coil type, the amplitude, phase (in case of harmonic analysis)
and the input function (in case of dynamic analyses) of the driving
current respectively voltage has to be given:

\begin{itemize}
\item Transient case
\begin{verbatim}
<value>       10           </value>
<dynamics>    timefile.dat </dynamics>
\end{verbatim}

\item Harmonic case
\begin{verbatim}
<value>       10           </value>
<phase>       90           </phase>
\end{verbatim}
\end{itemize}

\noindent
The name of the time function as well as the sepcification of the
phase is optional. If it's not given, a constant current with an
amplitude of {\bf value} is assumed.  
Voltage driven coils additionally require the total ohmic resistance
of the coil for solving the network equation:

\verb!<resistance>  10 </resistance>!

\subsubsection{2D Coils}
For all types of 2D-coils, an identification number is
necessary:

\verb! <id>  2  <id>!

\noindent
Due to the 2D- or axisymmetric modeling, only the cross-section of the
coil is modeled. The coil-id is
used to identify coil cross-sections belonging together (fed by the
same current). 
The sign of the id also specifies the direction of the current.


\subsubsection{3D Current Driven Coils}
For 3D current driven coils, the direction of the current has to be
given explicitly. In case of a straight current flow, three
possibilities can be given:

\verb!<currentFlow> direction </currentFlow>!

\noindent
with {\bf direction} standing for one of 
\begin{itemize}
\item xDir
\item yDir
\item zDir
\end{itemize}
   
For rotational current flow (e.g. in 3D-modeled axisymmetric coils),
additionally the midpoint and the direction of the rotational axis
must be defined:

\begin{verbatim}
<rotational>
   <midPointX> 0 </midPointX>
   <midPointY> 0 </midPointY>
   <midPointZ> 0 </midPointZ>
   <rotAxisX>   1 </rotAxisX>
   <rotAxisY>   0 </rotAxisY>
   <rotAxisZ>   0 </rotAxisZ>
</rotational>
\end{verbatim}

In this example, the midpoint lies at the origin, and the coil
rotation axis is equal to the x-axis.





\subsubsection{Coil Definition Examples}

The full description of a voltage driven coil in 2D with a winding
cross section of $10^{-3}\, \rm m^2$, a dynamic excitation of 10\,V with
the dynamic behaviour saved in the file "time.dat", 
and 10\,$\Omega$ resistance could be given as
\begin{verbatim}
<voltageCoil2d name="coil1">
   <windingCrossSection> 0.001  </windingCrossSection>
   <value>                  10  </value>
   <dynamics>         time.dat  </dynamics>
   <resistance>             10  </resistance>
   <id>                      2  </id>
</voltageCoil2d>
\end{verbatim}

A 3D rotational coil with the rotational axis parallel to the y-axis
including the point (2,2,0) should be defined as

\begin{verbatim}
<currentCoil3d name="coil2">
  <windingCrossSection> 1e-5 </windingCrossSection>
  <value>                100 </value>
  <dynamics>        time.dat </dynamics>
  <rotational>
    <midPointX>  2  </midPointX>
    <midPointY>  2  </midPointY>
    <midPointZ>  0  </midPointZ>
    <rotAxisX>    0  </rotAxisX>
    <rotAxisY>    1  </rotAxisY>
    <rotAxisZ>    0  </rotAxisZ>
  </rotational>
</currentCoil3d>
\end{verbatim}


\subsection{Permanent Magnets}
The use of permanent magnets in the simulation is very simple. In the
magnetic section, the command

\begin{verbatim}
<magnets>
   <magnet name="perm" orientX="0" orientY="0" orientZ="1" />
</magnets>
\end{verbatim}
has to list all domains, which hold a permanent magnet. The
orientation and strength of the permanent magnetic field is given by
the magnetization vector $\bM$ giben in Tesla described by the three
values {\bf orientX}, {\bf orientY} and {\bf orientZ}

\begin{equation}
  \label{eq:M}
  \bM = \left(
    \begin{tabular}{r}
      orientX\\
      orientY\\
      orientZ
    \end{tabular}
\right) \; .
\end{equation}

\subsubsection{Example}
{\footnotesize
\begin{verbatim}
<?xml version="1.0"?>
<cfsSimulation xmlns="http://www.cfs++.org">

   <analysis type="static"/>
   <geometry type="plane"/>
   <materialData file="mat.dat"/>

   <domain>
      <region name="myCoil" material="air"/>
      <region name="iron"   material="iron"/>
      <region name="air"    material="air"/>
      <nodes name="bound"/>
   </domain>

   <pdeList>
      <magnetic>

         <region name="myCoil" />
         <region name="iron"   />
         <region name="air"    />

         <bcsAndLoads>
            <dirichletHom   name="bound"/>
         </bcsAndLoads>

         <coils>
            <currentCoil2d name="myCoil">
               <!-- whole coil area: 100*5 mm^2 with 20 windings -->
               <windingCrossSection>50e-6 </windingCrossSection>
               <value>       1            </value>
               <dynamics>    time.dat     </dynamics>
               <id>          1            </id>
            </currentCoil2d>
         </coils>

         <storeResults>
            <elemResults type="magFluxDensity" region="myCoil"/>
            <elemResults type="magFluxDensity" region="iron"/>
            <elemResults type="magFluxDensity" region="air"/>
         </storeResults>

      </magnetic>
   </pdeList>
</cfsSimulation>
\end{verbatim}
}